% ============================================================
% APPENDIX D -- Computational Implementation
% Reorganised: Purpose -> Relationship -> Algorithms -> Solvers ->
% Complexity -> Software -> Reproducibility -> Limitations -> Future.
% ============================================================

\label{app:computation}

\begin{ataglance}[How to read this appendix]
This appendix is about arithmetic done by machines, and it is written so that
someone else could reproduce it or show it wrong.

Two separate pieces of software are described and it matters that they are not
confused. The empirical analysis of Chapters~\ref{ch:methods}
to~\ref{ch:results2} runs in R and is the work the thesis is judged on
(\S\ref{sec:D_relationship}). The dynamical simulation of
Appendices~\ref{app:dynamical} and~\ref{app:hemo_dynamics} is a proposed future
direction and is what the algorithms and solvers here describe
(\S\ref{sec:D_algorithms}--\S\ref{sec:D_solvers}).

The remaining sections say how long the computations would take
(\S\ref{sec:D_complexity}), what would have to be recorded for another
researcher to repeat them (\S\ref{sec:D_software}), and what the present
archive does not contain (\S\ref{sec:D_limitations}). That last section is the
one to read first if the question is what has actually been built.
\end{ataglance}

\section{Purpose}
\label{sec:D_purpose}
This appendix records the computational methods that implement the dynamical
extension of Appendices~\ref{app:dynamical} and \ref{app:hemo_dynamics}, so that
the proposed future modelling can be reproduced. The algorithms below are
schematic illustrations of the dynamical simulation, not the empirical analysis
pipeline.

\section{Relationship to the Empirical Analysis}
\label{sec:D_relationship}
The empirical measurement and prognostic models of Chapters~\ref{ch:methods} to
\ref{ch:results2} are intended to be estimated in R using \texttt{lavaan} for
structural equation modelling, \texttt{survival} for the Cox models, and
\texttt{mice} for imputation \citep{vanbuuren_2011}. The V2 archive includes an executable analysis
scaffold under \texttt{analysis/}; it validates inputs and records session
information, but cannot reproduce numerical results until a documented NHATS
extract and variable mapping are supplied. The dynamical simulation described
here is separate and remains a proposed future direction, not part of the
empirical validation.

The algorithms below therefore belong to the second of those two pieces of
software. They are written out in full because a specification that cannot be
checked is not a specification.

\section{Algorithms}
\label{sec:D_algorithms}
The simulation is organised in three layers that process biomedical, stochastic,
and policy inputs in turn.

\paragraph{Algorithm 1 (biomedical and DALY calculation).}
Disability weights and the DALY construction follow \citet{gbd2019}.
Given a disease profile, life expectancy $L_0$, disability weight $w$, and
survival function $S(t)$: compute the quality index $Q$; the severity-adjusted
life expectancy $L_{\text{adj}}=L_0 Q$; the life expectancy
$\mathbb{E}[L]=\int_0^\infty S(t)\,dt$; years lived with disability
$\mathrm{YLD}=w\,\mathbb{E}[L]$; years of life lost $\mathrm{YLL}=L_0-\mathbb{E}[L]$;
$\mathrm{DALY}=\mathrm{YLD}+\mathrm{YLL}$; and biomedical potential
$BP=(L_{\text{adj}}/L_0)(1-\mathrm{DALY}/L_0)$.
The identity $\mathrm{DALY}=\mathrm{YLD}+\mathrm{YLL}$ follows directly from
the definitions above. A future reference implementation should test both this
identity and the numerical quadrature for $\mathbb{E}[L]$.

\begin{reading}[What Algorithm 1 computes, in clinical terms]
The algorithm turns a disease profile into a single number for biomedical
potential, and every step of it is standard burden-of-disease accounting.

Years of life lost is how much shorter life is expected to be than it would
otherwise have been. Years lived with disability is the remaining life weighted
by how impaired it is. Their sum is the disability-adjusted life year, which is
the currency the Global Burden of Disease study reports in.

For Mr~B the two components separate cleanly. His years of life lost are now
modest, because treatment has largely closed the survival gap that defined
haemophilia before the 1990s. His years lived with disability are substantial,
because a destroyed knee and chronic pain carry a high disability weight
for the whole of his remaining life. A measure that reported only mortality
would record him as almost unaffected, which is the reason the DALY was
invented and the reason biomedical potential is built from both terms.

The final line, $BP=(L_{\text{adj}}/L_0)(1-\mathrm{DALY}/L_0)$, rescales that
onto the unit interval the rest of the framework uses. It is a normalisation
step and carries no separate claim.
\end{reading}

\begin{workedexample}[apd-daly]{Mr~B's burden, split into its two parts}
Suppose a man of Mr~B's age would without haemophilia expect $L_0=30$ further
years, that treatment has brought his life expectancy to $\mathbb E[L]=28$
years, and that his joint damage and chronic pain carry a disability weight
$w=0.28$. Take the quality index as $Q=0.72$.

\begin{wecode}[Python]
from hpv.economics import daly, biomedical_potential

yld, yll, total = daly(L0=30.0, w=0.28, E_L=28.0)
print(f"YLD {yld:5.2f}   YLL {yll:5.2f}   DALY {total:5.2f}")
print(f"biomedical potential {biomedical_potential(30.0, 0.72, total):.4f}")

# the same man, in the era before prophylaxis reached him
yld0, yll0, total0 = daly(L0=30.0, w=0.28, E_L=15.0)
print(f"\npre-prophylaxis: YLD {yld0:5.2f}   YLL {yll0:5.2f}   "
      f"DALY {total0:5.2f}")
\end{wecode}

\noindent Today: $\mathrm{YLD}=7.84$, $\mathrm{YLL}=2.00$, so a DALY of $9.84$
of which eighty per cent is disability. In the era before prophylaxis reached
him the same weights give $\mathrm{YLD}=4.20$ and $\mathrm{YLL}=15.00$, a DALY
of $19.20$ of which seventy-eight per cent is death.

The total halved and its composition inverted. That is the whole argument for
this thesis in one calculation: a measure that tracked mortality alone would
record the modern Mr~B as almost unaffected, having correctly recorded the
1970s Mr~B as gravely affected, and would have nothing to say about the twenty
years of reduced mobility that now constitute his burden.

\weanswer{$\mathrm{DALY}=9.84$, eighty per cent of it disability, giving
$BP=0.484$. Asserted by \texttt{test\_algorithm\_one\_daly\_identity} and
\texttt{test\_worked\_example\_mr\_b\_daly\_split}.}
\end{workedexample}

The cost side of the same decision needs its own arithmetic, and the difficulty
there is not the point estimate but the interval around it.

\paragraph{Algorithm 2 (Monte Carlo cost-effectiveness).}
Draw costs and effects for two strategies, form $\Delta C$, $\Delta E$, and
$\mathrm{ICER}=\Delta C/\Delta E$; report the ratio-of-means ICER, a bootstrap
\citep{efron_tibshirani_1993} confidence interval, and the incremental net
monetary benefit at a willingness-to-pay threshold \citep{briggs_2006}. Recovery of a planted ICER with a bracketing
interval should be verified by a unit test before a numerical result is reported.

\begin{reading}[Why an ICER needs a bootstrap and not a formula]
The incremental cost-effectiveness ratio is a ratio of two uncertain
quantities, and ratios misbehave. When the denominator $\Delta E$ is small and
its confidence interval crosses zero, the ratio's interval is not an interval
at all: it covers everything except a finite middle region, and the usual
standard-error formula returns a number that means nothing.

The clinical intuition is the same one behind a relative risk computed on three
events. The point estimate looks precise and the ratio is unstable because the
denominator is nearly nothing.

This is why the net monetary benefit is reported alongside. Rearranging
$\lambda\Delta E-\Delta C$ removes the division and gives a quantity whose
uncertainty behaves, at the cost of requiring the threshold $\lambda$ to be
named rather than left implicit in a ratio. Naming it is an improvement, not a
concession.
\end{reading}

Both algorithms assume the underlying differential equations can be integrated
accurately. That is a property of the solver rather than of the model, and it
has to be demonstrated rather than assumed.

\section{Numerical and Stochastic Solvers}
\label{sec:D_solvers}
\begin{itemize}
  \item \textbf{Deterministic.} Explicit Runge--Kutta 4(5) \citep{butcher_2008} for ODEs, with stiff
  variants (Radau, BDF) invoked adaptively; a fixed-step RK4 is used where the
  step count must be controlled.
  \item \textbf{Stochastic.} Euler--Maruyama and Milstein schemes \citep{kloeden_platen_1992} for
  $dX=f\,dt+g\,dW_t$; the Milstein correction raises the strong order to one.
  \item \textbf{Population diffusion PDE.} Conservative finite-volume assembly
  with an implicit (backward-Euler / Crank--Nicolson) time step; alternating-
  direction implicit schemes generalise this to higher dimensions.
\end{itemize}
For a sufficiently smooth problem in the asymptotic regime, fourth-order RK4
should show an error ratio approaching $16$ under step halving; this is a
recommended verification test rather than a result established in this archive.

\begin{reading}[What the order-of-convergence test actually checks]
Halve the step size and the error of a method of order $q$ should fall by a
factor of $2^q$: sixteen for a fourth-order scheme, thirty-two for a
fifth-order one. The embedded Dormand--Prince pair used here advances at fifth
order and carries a fourth-order estimate for step control, so the ratio to
look for depends on which solution is propagated, and the test has to state
which. Observing the predicted ratio is the standard way to show that a solver
is doing what it claims.

It is worth being clear about what the test does and does not establish. It
checks the arithmetic, not the biology. A solver can integrate a wrong equation
to machine precision, and the convergence test will pass. Passing it means the
numbers in a figure are the numbers the model implies; it does not mean the
model implies anything true about a patient.

The distinction matters here because the dynamical appendices make no empirical
claims. Their whole content is that certain equations have certain properties,
and a solver that fails this test would leave even that unestablished.
\end{reading}

\begin{workedexample}[apd-order]{Running the convergence test rather than asserting it}
The claim is that halving the step cuts a fourth-order method's error by
sixteen. The way to establish it is to run it. Integrate the deterministic
model of Appendix~\ref{app:dynamical} from a fixed start over four time units
with fixed-step RK4, taking a very fine run as the reference.

\begin{wecode}[Python]
from hpv.dynamics import integrate_ode

P = dict(Vmax=100.0, Km=10.0, T=20.0, k_inh=0.05, Y=2.0,
         gamma=30.0, Kd=50.0, n=2.0, Z=80.0, k_elim=0.8)
ref = integrate_ode(0.0, 4.0, P, dt=1e-5)

prev = None
for dt in (0.4, 0.2, 0.1, 0.05):
    err = abs(integrate_ode(0.0, 4.0, P, dt=dt) - ref)
    ratio = "" if prev is None else f"   ratio {prev / err:5.2f}"
    print(f"dt {dt:<5} error {err:.3e}{ratio}")
    prev = err
\end{wecode}

\noindent Errors of $1.83\times10^{-3}$, $9.81\times10^{-5}$,
$5.68\times10^{-6}$ and $3.42\times10^{-7}$, giving ratios of $18.61$, $17.25$
and $16.61$. The ratios approach sixteen from above as the step shrinks, which
is what fourth-order convergence looks like in practice: the asymptotic rate is
reached from the direction of the higher-order terms that dominate at coarse
steps.

Two cautions belong with this. A ratio near sixteen shows the integrator is
solving the equation it was given to the order it claims; it says nothing about
whether that equation describes a patient. And the reference run is itself
numerical, so the last row is measuring the coarse run against a fine run
rather than against the truth; the test stops refining before that distinction
starts to matter.

\weanswer{Error ratios of $18.61$, $17.25$ and $16.61$ under successive step
halving, approaching the fourth-order value of $16$ from above. Asserted by
\texttt{test\_worked\_example\_apd\_order\_of\_convergence}.}
\end{workedexample}

Accuracy settled, the remaining practical question is whether the computations
finish in reasonable time.

\begin{lstlisting}[language=Python,basicstyle=\footnotesize\ttfamily,frame=single]
import numpy as np
from scipy.integrate import solve_ivp

def disease_model(t, y, params):
    X = y[0]
    Vmax, Km, k_inh, Kd, n, k_elim, Y, Z, T = params
    f = Vmax * T / (Km + T)          # Michaelis-Menten treatment
    g = k_inh * Y * X                # bilinear inhibition
    h = Z**n / (Kd**n + Z**n)        # Hill stabilisation
    return [f - g + h - k_elim * X]

sol = solve_ivp(disease_model, [0, 100], [1.0], args=(params,), method='RK45')
\end{lstlisting}

\section{Computational Complexity}
\label{sec:D_complexity}
The cost of each solver scales as follows, with $n$ time steps, $N$ paths, $M$
Monte Carlo draws, and $(n_x,n_t)$ the space and time resolution of the PDE.
\begin{center}
\begin{tabular}{ll}
\toprule
Component & Complexity \\
\midrule
ODE (RK4, 4 evaluations per step) & $O(n)$ \\
SDE (Euler--Maruyama, $N$ paths)  & $O(nN)$ \\
Monte Carlo ($M$ draws)           & $O(MN)$ \\
PDE (Fokker--Planck, sparse solve)& $O(n_x n_t)$ \\
\bottomrule
\end{tabular}
\end{center}
For a fixed-step implementation, RK4 uses $4n$ right-hand-side evaluations and
Euler--Maruyama updates $nN$ path--time cells. These counts should be asserted in
unit tests when the future simulation package is implemented.

The table has one line worth noticing. The single-patient ODE costs $O(n)$ and
the population Fokker--Planck solve costs $O(n_xn_t)$, so following a whole
cohort as a density is cheaper than simulating enough individual paths to
estimate the same distribution. That is the practical argument for the
population layer of Appendix~\ref{app:dynamical}, alongside the conceptual one.

What none of these costs include is the work of making a computation
reproducible by someone else, which is where most of the effort in practice
goes.

\section{Software and Reproducibility}
\label{sec:D_software}
The proposed dynamical implementation would use NumPy for arrays, SciPy for
integration and linear algebra, and optionally Numba for performance-critical
loops. Reproducibility requirements are a fixed random seed, a locked software
environment, machine-readable parameters, unit tests for limiting cases and
convergence, and scripts that regenerate every reported figure. These are design
requirements, not claims about software included in the present archive.

That last sentence is the honest position, and the next section states its
consequences plainly rather than leaving them to be inferred.

\section{Limitations and Future Implementation}
\label{sec:D_limitations}
The code is a reference implementation of a proposed modelling direction, not a
clinical tool: parameters are illustrative, no patient data is used, and the
population layer is one-dimensional. Future work would calibrate parameters to
longitudinal factor-level data, extend the PDE to several biomarkers, and couple
the dynamical state to the validated latent score through the time-varying form
of Appendix~\ref{app:dynamical}.

\section{Summary}
\label{sec:D_summary}
This appendix collected the algorithms, solvers, complexity, and reproducibility
requirements for the dynamical models of Appendices~\ref{app:dynamical} and
\ref{app:hemo_dynamics}. It is deliberately separate from the empirical R
pipeline of the main text. No claim of machine verification is made until the
future simulation implementation and its tests are deposited with the thesis.
